%% This LaTeX-file was created by <imt> Fri May  8 13:10:31 1998
%% LyX 0.12 (C) 1995-1998 by Matthias Ettrich and the LyX Team

%% Do not edit this file unless you know what you are doing.
\documentclass{article}
\usepackage[T1]{fontenc}
\usepackage{makeidx}
\makeindex
\usepackage{graphics}

\makeatletter


%%%%%%%%%%%%%%%%%%%%%%%%%%%%%% LyX specific LaTeX commands.
\newcommand{\LyX}{L\kern-.1667em\lower.25em\hbox{Y}\kern-.125emX\spacefactor1000}
\newcommand{\noun}[1]{\textsc{#1}}

\makeatother

\begin{document}


\title{Management Report }


\author{System Design Group Seven}

\maketitle
\vspace{0.3cm}
{\centering \includegraphics{../web/images/EUcrest-official-noborder.ps} \par}
\vspace{0.3cm}

\newpage
\tableofcontents

\printindex


\newpage
\section{April 21 (11am - 12am)}

In our initial meeting it was decided to concentrate on accumulating ideas for
the product and the group. We had to decide on the roles each person would be
taking and how we would begin to tackle the design specification. Each person
stated their preferred interests and the group was split up accordingly

\begin{itemize}
\item Ian Clarke (Computer Science and Artificial Intelligence) Software 
\item Stephen Cook (Computer Science and Management) Group coordinator and hardware
\item Stuart Cox (Computer Science and Physics) Hardware
\item Duncan Irvine (Computer Science) Software
\item Stephen Kitt (Computer Science and Artificial Intelligence) Software
\item Edward Knowelden (Computer Science andArtificial Intelligence) Software
\item James Northrop (Computer Science) Hardware 
\item Iain Todd (Computer Science) Feasibility study and documentation
\item Gordon Brebner Group Ssupervisor
\end{itemize}
This enabled us to have three groups each working on specific parts of the project.
The three groups will meet up with the group supervisor three times a week.
Eacg group will conuct there own meetings wehnenver that may may bo appropriate.
Each of the three groups will have to keep in daily contact with each other,
so that their efforts can be coordinated towards one goal.

We drew up a time chart for the next four weeks, duly noting the three Deadlines
that were imposed on us, and adding our own deadlines. These deadlines included
when the design should be finalized, the deadline for hardware and the presentation.
This chart will be amended as the project progresses, and our goals became clearer.
\noun{<see appendix 1> }


\subsection{General Ideas}

Our brainstorming produced specific ideas for the different areas of the practical.


\subsubsection{Security}

Swipe cards would be cheap and simple to implement. They also seem to be the
favored mechanism for security eg the North Machine Halls, The AI department.
Dallas rings are as portable as swipe cards, but we have little information
about them. Smart cards would allow us to provide information on the cards,
but we are unsure about the cost. We would also need to decide wether a backup
of a PIN number would be required by each method


\subsubsection{Special Features}

Features such as voice, face or iris recognition were considered but there is
a worry over the likliehood of success. We are unsure wether any of these methods
have been perfected. If we were to implement one of these features we would
need to make sure we had a contignancy plan.


\subsubsection{Market}

The initial specification states that the system will be designed as a replacement
for the standard door-entry system of residential flats. We would need to investigate
wether residents would accept this product and if they didn't we would investigate
new markets.


\subsection{Goals}

The task of investigating these ideas was split up between the group. We set
a deadline of Thursday Wednesday April 22 to make decisions on our market and
our features. The features of the unit will be dependent on the market that
we aim for.


\newpage
\section{April 22 (10am - 11am)}


\subsection{Research Results}


\subsubsection{Security }

Smart cards are expensive but there is an option of making our own using BT
phone cards. The principle is quite simple but the implementation may not be.
We can not buy smart cards as they are expensive, the card writer alone would
cost more than our budget. Dallas rings are not the cheapest option, but we
would probably be the only group using them.


\subsubsection{Features }

Examining the area of image recognition as a security mechanism; each user would
be given a card with a barcode on it. The user would show this to the camera
(which would recognize the code) and then type in a PIN, if the pin and code
were correct then they would gain entry.We consider the system to be feasible,
the algorithm to recognize the image is quite simple. Advantageously the cost
would be low, the image would only be printed onto a piece of laminate, no special
equipment would have to be bought. The limitation to this method is that we
have not been able to test the camera to check the quality of the image. The
image will need to be quite clear to make out the bar-code.


\subsubsection{Market}

Our initial research shows apathy from most residents, the standard response
to the idea is why should they spend money on a new system when they allready
have one that meets all their needs. It was decided that we should target the
commercial buyer. We had the specific idea of targeting school security. In
light of Dunblane there should be an increased market for low cost security.


\subsection{Reports}

The hardware group were looking into the chips on the board and the connection
methods. They will need time to get familiarized with the board and the chips
involved.

It was suggested by supervisor that we were expanding a lot of time on the technical
details of the product while neglecting some basic ideas. How would the system
cope with general situations; switching it off, the machine been dropped and
the PC not been switched on. 


\subsubsection{Deadlines}

Documentation would start on the Design Specification today, so that we could
meet our first Deadline. We would also make a final system design decision at
a meeting later in the day.


\newpage
\section{April 22 (4pm - 5pm)}


\subsection{Final Research Results}

We have fully Investigated Dallas rings, smart cards and swipe cards. Dallas
rings cost \pounds{}2 each, cheaper when bought in bulk. The reader is \pounds{}5.
The smart card will not be a good option as there is no supplier of blank cards.
There is only a small amount of literature available, finding information about
the cards is proving difficult. If we were to use this method we would be carrying
out a lot of work into an area were we may get little help. Swipe cards readers
cost \pounds{}25 but the cards themselves only have a lifespan of one to two
years. We have sent away for more data about the rings and the swipe cards.
The principal problem with the rings are that the readers is not vandal proof.
This would not be a great problem as it is a protype we are designing. The problem
with swipe cards is that the programmer is too expensive.

Continuing our investigation into image recognition, we think the good points
would be: 

\begin{itemize}
\item It looks impressive, 
\item It would be very convenient to use 
\item It would be cheap. 
\end{itemize}
The downside is that we would be unable to test it until near the end of the
project. We would have to complete most of the software and the hardware without
the guarantee of a working product. Until we get a picture from the Quick-Cam
then we will not know wether it is feasible. 


\subsection*{Reports}


\subsection{Hardware}

The Hardware team have finished their investigation into the chips and the board.
They have booked consultation time to help them start up, and their initial
thought is that the hardware is going to be more difficult than first thought. 


\subsection{Software}

The software team has checked out the audio capabilities on the Linux machine
and it seems as if it will operate correctly. The video capabilities seems to
be about one frame per second. They will continue to examine the serial connection
between the Motorola board and the PC.


\subsection{Documentation}

No company deals directly in school security. Most firms try to make there product
available to any market that desires access control. The case of babies been
snatched form maternity wards highlights the needs of hospitals for security.


\subsection{Decisions}

We made the decision to use Dallas Rings. We will not fully rule out image recognition
until we get the camera operating. Our target market will be public institutions
such as schools and hospitals.


\subsection{Plan}

Documentation will complete the Requirement spec for 12pm April 23. Then it
can be corrected and updated by the rest of the group. Hardware use some consultation
time and software will have an early morning meeting to allocate tasks amongst
themselves.


\newpage
\section{April 24 (10.30pm - 11pm)}


\subsection*{Reports}

The camera was set up this morning and the quality of the picture is not clear
enough to continue with image recognition. We tested holding up a card with
a barcode to the camera and it was very unclear the gaps between the black and
the white bands. We could now dedicate our time fully to the Dallas Rings.


\subsection{Hardware }

They have looked at the capabilities of the chips and have started work on timing
diagrams. Once they have the design they should have no problems with the construction.
They aim to have the hardware up and running by the start of the third week.
To achieve this they will order some chips immediately so that they will not
be held back later in the project.


\subsection{Software }

They are considering the software requirements. The software will be split into
modules and consideration will have to be given about how the two main areas
will connect. The software team have split up to consider their individual assignments
but they hope to have something working by Tuesday. They have considered all
the features they wish to include (see Design Specification)


\subsection{Documentation }

Has to meet the deadline of 4.30 on April 23. Will check the documentation with
Stuart Anderson as soon as possible, to correct the structure and contents.
Our market research suggests that schools should be our target market. We have
contacted schools to gauge their opinion on the unit, and we hope to get the
results of our residents survey over the weekend. 

In light of the Cullen inquiry the security should be should be good enough
to allay parents worries but not be to restrictive to an environment that should
be welcoming. The system will be tailored towards their needs and will have
to be configurable. Different schools will have competing needs. The research
indicates a market with the NHS; especially maternity wards which increasingly
need a low level of security and identification. The advantage of the Dallas
ring over swipe cards is that they will not be wiped in a highly magnetized
environment such as a hospital.


\subsection{Plan}

To set a clear deadline for the Hardware group, as the software has a limited
dependency on their finished project. If the worst comes to the worst they should
investigate buying the design. Any any other details are covered by the time
chart.


\newpage
\section{April 27 (10.30am - 11am) }


\subsection*{Reports}


\subsection{Hardware}

It is not possible to buy the hardware design but you can get a good idea on
how to go about it by using up consultation time. They hope to start construction
by the end of the week and have something working for Monday. They will order
some of the hardware today.


\subsection{Software}

The video is almost complete but it needs to run faster. They consider progress
to be good. The Audio specs has now being decided and they are investigating
the interface to see how it will multi-task. They have an idea of using the
SSCOP code for the serial link, but this may be overcomplicating the situation.


\subsection{Documentation}

Any hand-ins should have a system of proof reading. This should be done by a
designated member of the group with good english skills. Hand-ins should be
written in a style that will assume no prior knowledge of the subject. The feasibility
study is still on going.


\subsection{Plan}

The Software team will continue to work on the serial cable and should consider
booking consultation time if they have any trouble with it. Programming the
video software continues and they will delegate other tasks amongst themselves.
The hardware team are making progress, and have ordered the parts. Documentation
should start on the feasibility study and there shuld be corrections made to
the design specification


\newpage
\section{April 29 (3pm - 4pm)}


\subsection*{Reports}


\subsection{Hardware}

They have ordered and received some of the parts, but they will need to get
a few others. They are nearly ready to start construction; Tim Hopkins has checked
their design and seems to have no problem with it. They have started some rudimentary
testing. For the Feasibility study it should be noted that two of the chips
we are using are obsolete.


\subsection{Software}

There is a snag with the interface, we may not be able to get the audio and
video working together. The camera will record pictures until the buzzer sounds
and then only the audio will be active. The user could then check the previous
video frames for a clear picture. We should aim to have the simpler design running
before we expand it. It is noted that this design does not meet our minimum
spec.

The video compression could run at one frame per second. Once all the hardware
is finished we we will be able to test the client side software. The audio has
has a choice of two different standards one of which might be to hardware intensive
to proceed with.


\subsection{Documentation}

Has made contact with our target groups to gauge a response to our system. Continues
to work on the Feasibility study. Should contact the Royal Infirmary or the
Scottish Schools Research Center. A phone can be obtained from the secretaries
office.


\subsection{Plan}

See time chart


\newpage
\section{April 30 (3am - 4am) }


\subsection*{Reports}


\subsection{Hardware}

They have completed most of the design. The audio section is proving problematic;
it is likely that the hand-in tommorow will contain a faulty design. They hope
to have finished construction by Sunday. Some Other Groups are just handing
in block diagrams; but we will be handing in documents that contain as much
detail as is possible. Comprising of a technical description, flow of data,
timing diagrams etc.


\subsection{Software}

They have designed a state diagram, and have member of the group has been allocated
specific parts to write up and implement. They will write up a description,
flow of data and the transitions for each state. The GUI is designed but they
wil need to work out how the client side operates. The audio will have something
up and running by Tuesday. Integration of the states can only take place once
all the individual states have been finalized.


\subsection{Documentation}

The Feasibility study is still on going.


\subsection{Plan}

The individual parts of the hardware and software documentation have to be written
by 9am 1st May. Documentation will then put all the parts together for 10.30am.
The document will be proof read until approximately 12am to have it completed
for 1pm


\newpage
\section{May 1st (3am - 4am) }


\subsection*{Reports}


\subsection{Hardware}

The hardware is theoretically correct. They have checked it with Tim and hope
to start building it. The construction will take place in stages as they need
to complete some testing of chips.


\subsection{Software}

See the state diagram and the previous minutes. They have allocated 2-3 days
next week for integration. Each state has now been documented by each individual,
and the document is been produced. GB wants us to write about the functionality
rather documenting all the code. Their only worries are with the Audio implementation.
It is not known how the Motorola will cope with the code. Last years groups
had a lot of trouble porting their code over. 


\subsection{Documentation}

We have contacted the SSERC and have managed to get some quotes from them. The
document itself is nearly complete. It will be worked on by the Hardware and
Software team to clear it of any inaccuracies. Work will now start on updating
the minutes into a readable management report.


\subsection{Plan}

see time chart


\newpage
\section{May 4th (3am - 4am)}


\subsection*{Reports}


\subsection{Hardware}

The hardware Is looking good everything has been tested and it is just a matter
of putting it all together. They have built everything apart from the audio
compression.


\subsection{Software}

Video is as far on as it can get, without the serial being implemented. The
video may take up to a day to test once the serial is complete. The serial is
continuing. Work is nearly complete on the Linux side, and it is theoretically
correct on the Motorola side. They are integrating the video with the GUI. They
have an audio compression program that is being read through and optimized.


\subsection{Documentation}

The Superrvisor was surprised at the level of detail in the technical document
but praised the hand in, stating he had confidence that we would manage to complete
our design. Minor points were that the figures did not always relate to diagrams;
and detail was sketchier in some place than others. It should be made clearer
that the mechanical lock is a feature of the door and not the unit. We should
also state more clearly how refuse entry works.

The Feasibility study is still on going. It is now at a point were it should
be handed over to somebody else, so they can contribute a different outlook.
Progress should now be made on the presentation and the market research.


\subsection{Plan}

Documentation to start on presentation. Towards the end of the week the hardware
team will split up into software and documentation


\newpage
\section{May 7th (3am - 4am)}


\subsection*{Reports}


\subsection{Hardware}

Everything has been built, wired and individually tested. They need to test
and run the system as one unit; but are hopeful of meeting our self imposed
Friday deadline. In case anything goes wrong we have copies of all the chips
and will be able to replace any broken part. One slight problem is that the
video interferes with serial and vice versa, but this is being examined.


\subsection{Software}

The database has been tested on the Motorola and seems to be working correctly
The testing will continue for the rest of the day. When we get more Dallas rings
then we will be able to use correct values. The GUI is still in development
but it shouldn't prove a problem if everything else is up and running. The serial
is correct in one direction and work has started on the other way. Video has
been tested with the serial and is proving slightly buggy. The first image transmitted
is correct the rest are corrupted. The audio standard has been changed from
ADPCM to a more simpler method. This comes about as ADPCM was proving to be
to hardware intensive for our hardware.


\subsection{Documentation}

Worked on 3D packages for the presentation, conclusion is that 3D studio max
will have to be used. Has finished the management report up until present.


\subsection{Plan}

The hardware team will be splitting up into marketing and documentation, one
of them well liase with the software team to explain any design points. Everyone
is confident of meeting their individual deadlines for this Friday. The deadline
envolves having a working unit to demonstrate to our group supervisor.

\end{document}
